
On
sait que pour tout schéma régulier ayant un faisceau ample, les groupes de Grothendieck formés avec les faisceaux localement libres de type fini et les faisceaux cohérents sont identiques. C'est donc en particulier le cas pour \(Y'\).

Soit \(\underline{E}\) un \(\underline{O}_{X'}\)-Module cohérent à support dans \(Y'\); alors \(\underline{E}\) est un \(\underline{O}_{X'}\)-Module parfait. En effet, il l'est aux points de \(Y'\), puisque \(X'\) est régulier en ces points, et aux points de \(X'-Y'\) puisqu'il y est nul.

Si \(E\) est la classe de \(\underline{E}\) dans \(K^{\bullet}(X')\), \(E\) appartient à l'image de \(j_*\). En effet, on peut filtrer \(\underline{E}\) par les \(\underline{J}'^n.\underline{E}\). Cette filtration est finie, car, \(X'\) étant noethérien, il existe un entier \(n_0\) tel que \(\underline{J}'^{\,n_0}.\underline{E}=0\); les quotients de cette filtration sont alors des \(\underline{O}_{Y'}\)-Modules, donc leur classe dans \(K^{\bullet}(X')\) appartient à l'image de \(j_*\). Comme \(E\) est la somme des classes de ces quotients, \(E\) y appartient aussi.

Soit maintenant \(x\in K^{\bullet}(X')\). On peut trouver un complexe parfait \(\underline{M}^\bullet\) dont la classe soit \(x\). Considérons le triangle distingué:
\[
\begin{tikzcd}[column sep=large,row sep=large]
& \underline{N}^\bullet \arrow[dl] &\\
\underline{M}^\bullet \arrow[rr] && \underline{\underline{L}}f^*(\underline{\underline{R}}f_*(\underline{M}^\bullet)) \arrow[ul] ,
\end{tikzcd}
\]
la flèche horizontale étant l'homomorphisme canonique. Cette flèche est un isomorphisme sur le complémentaire de \(Y'\), et par conséquent les faisceaux de cohomologie de \(\underline{N}^\bullet\) sont des faisceaux cohérents à support dans \(Y'\). D'après ce qui précède, ce sont des \(\underline{O}_{X'}\)-Modules parfaits, et leur classe appartient à l'image de \(j_*\). On a donc:
\[
x-f^*(f_*(x))\in\operatorname{Im}(j_*)\,;
\]
par suite, l'homomorphisme \(v\) est surjectif.



% Macros for Exposé VII continuation
\providecommand{\Q}{\mathbb Q}
\providecommand{\bfP}{\mathbf P}
\providecommand{\gr}{\operatorname{gr}}
\providecommand{\Id}{\operatorname{Id}}
\providecommand{\Fil}{\operatorname{Fil}}
\providecommand{\Imf}{\operatorname{Im}}
\providecommand{\rank}{\operatorname{rank}}
\providecommand{\codim}{\operatorname{codim}}
\providecommand{\cJ}{\mathcal J}
\providecommand{\cI}{\mathcal I}
\providecommand{\cN}{\mathcal N}
\providecommand{\cE}{\mathcal E}
\providecommand{\cO}{\mathcal O}


